\section{Sceau dodécaphasique et lecture rationnelle réversible}
\label{sec:sello}

Le registre terminal possède une coordonnée visible et une coordonnée
signée. Le sceau s'obtient à partir de cette dernière par la transformation
de Hadamard par orbites ; une lecture scalaire ultérieure permet de retrouver
le sceau lorsque son cadre de représentation est conservé. Cette propriété
précise ce que signifie ici une définition structurelle d'une constante :
les applications de production et de récupération sont explicitées, en plus de sa valeur.

Soient $U\in\Z^{12}$ la coordonnée signée produite par le registre,
$H_4$ la matrice de Hadamard développée à la section~\ref{sec:action}
et $P_{\rm orb}$ la permutation qui regroupe
\[
 (1,4,7,10),\qquad(2,5,8,11),\qquad(3,6,9,12).
\]
Définissons
\[
 \mathcal H_{12}=P_{\rm orb}^{-1}
 (H_4\oplus H_4\oplus H_4)P_{\rm orb}.
\]
On a $\mathcal H_{12}^2=4I$. Lorsque
$\mathcal H_{12}U\equiv0\pmod4$, le sceau entier est
\[
 K=\tfrac14\mathcal H_{12}U,
 \qquad U=\mathcal H_{12}K.
\]
Les différences transversales $D_3K,D_4K$ et la charge $Q(K)$ retrouvent
les canaux spécifiés par le registre (voir \ref{prop:k-transversal}).
Le sceau est ainsi une coordonnée reconstructible de ce registre orienté ;
la profondeur et l’histoire complète restent des champs distincts.

\begin{proposition}[Récupération des blocs et transport de l'origine]
Une base entière $b\ge2$ et douze positions ordonnées étant fixées, soit
$K_j\in\{0,\ldots,b-1\}$. Définissons
\[
 I(K)=\sum_{j=1}^{12}K_jb^{12-j},\qquad
 \kappa(K)=\frac{I(K)}{b^{12}-1}.
\]
La lecture exacte $\kappa$ détermine de manière unique les douze blocs.
Pour la rotation $\rho K=(K_2,\ldots,K_{12},K_1)$,
\[
 \boxed{\kappa(\rho K)=b\kappa(K)-K_1.}
\]
\end{proposition}
\begin{proof}
L'entier $I=(b^{12}-1)\kappa$ se décompose de manière unique en douze
positions de base $b$, en complétant par des zéros les positions initiales
nécessaires. D'autre part,
\[
 I(\rho K)=bI(K)-K_1(b^{12}-1).
\]
Diviser par $b^{12}-1$ donne la loi de transformation. Pour le bloc
constant $K_j=b-1$, la lecture est $\kappa=1$ ; le type de représentation
périodique fixé conserve également sa récupération.
\end{proof}

Dans la base $1000$ du registre, les blocs de trois chiffres demeurent
récupérables, y compris les zéros initiaux. La lecture rationnelle
périodique conserve davantage qu'une approximation décimale tronquée.
Lorsqu'on modifie l'origine de phase, sa valeur change généralement selon
l'identité précédente ; au bout de douze rotations, elle revient.
L'avancée de mémoire de l'état TPK est une transformation supplémentaire :
la périodicité du lecteur n'identifie pas les histoires accumulées.

\begin{remark}[Complexité, entropie et transport]
La complexité des mots et leur transport par la retenue vers
$e+\pi$ sont des propriétés différentes de l’entropie d’une dynamique
symbolique, de l’entropie de l’état réduit et de la production thermodynamique.
Chacune nécessite sa définition et une application de transport spécifique.
La réversibilité du registre démontrée dans cette section conserve
les blocs et leur origine de phase ; à elle seule, elle n’établit pas de
classification arithmétique de $e+\pi$ et n’identifie pas les trois notions
d’entropie.
\end{remark}
